\section{Finite-cross-section conductor discretization and ports}
\label{sec:conductor-discretization}

This section fixes the mathematical meaning of the conductor teacher so that an
implementation does not reintroduce a hidden world mesh.

\subsection{Intrinsic coordinates and modal current space}

Let $s\in[0,L]$ denote arc length along a conductor centerline and
$\xi=(\xi_1,\xi_2)\in\Sigma(s)$ the transported cross-section coordinate.
Using the local frame
\[
  \{e_t(s),e_1(s),e_2(s)\},
\]
write
\[
 \mathbf J(s,\xi)
 =
 e_t(s)J_t(s,\xi)
 +e_1(s)J_1(s,\xi)
 +e_2(s)J_2(s,\xi).
\]
The longitudinal and transverse components are expanded as
\[
 J_t(s,\xi)
 =
 \sum_{\ell=1}^{N_s}\sum_{m=1}^{N_t}
 a_{\ell m}\chi_\ell(s)\psi_m^{(t)}(\xi;s),
\]
\[
 \mathbf J_\perp(s,\xi)
 =
 \sum_{\ell=1}^{N_s}\sum_{m=1}^{N_\perp}
 b_{\ell m}\chi_\ell(s)\bm\psi_m^{(\perp)}(\xi;s).
\]
The cross-section modes are defined on a reference superellipse and mapped to
the physical cross-section. They may be orthonormalized in the conductive
energy inner product
\[
 \inner{u}{v}_{\sigma^{-1}}
 =
 \int_{\Sigma}
 u\adj\sigma^{-1}v\,\dd A .
\]
At minimum the basis contains a uniform longitudinal mode and a sequence of
boundary-layer/edge modes able to represent skin and proximity concentration.

\subsection{Charge space and exact discrete continuity}

Choose a scalar charge basis
\[
  \rho(s,\xi)
  =
  \sum_{r=1}^{N_\rho} q_r\nu_r(s,\xi).
\]
Let $D$ be the discrete divergence map from current coefficients
$c=[a^\T,b^\T]^\T$ to charge-test coefficients:
\[
  (Dc)_r
  =
  \int_{\Omega_c}
  \nu_r^\ast\nabla\cdot\mathbf J(c)\,\dd V .
\]
Let $M_\rho$ be the charge mass matrix. For the \emph{total} conductor
current, including its terminal flux, harmonic charge continuity is imposed
algebraically,
\[
  Dc+j\omega M_\rho q=0.
  \label{eq:discrete-continuity}
\]
The terminal current is imposed independently by the exact port constraint
$\mathcal I c=i$ in Section~\ref{sec:port-schur}. Continuity is therefore not
a penalty term and cannot drift during iterative solves.

If an implementation introduces a compatible current lifting
\[
  c=c_0+B_I i,\qquad \mathcal I B_I=I,
\]
then the zero-port-current correction satisfies
\[
  Dc_0+j\omega M_\rho q
  =S_\rho i,
  \qquad
  S_\rho=-DB_I.
\]
Thus a nonzero $S_\rho i$ is a derived lifting source, not an additional
physical injection to be combined with the total-current port constraint.
Summing Equation~\eqref{eq:discrete-continuity} over a conductor gives exact
net-charge compatibility with its terminal fluxes.

\subsection{Port definition}

Each electrical port $p$ is a pair of oriented terminal regions
$(\Gamma_p^+,\Gamma_p^-)$. Define terminal current functionals
\[
 I_p(\mathbf J)
 =
 \int_{\Gamma_p^+}\mathbf J\cdot n\,\dd A
 =
 -\int_{\Gamma_p^-}\mathbf J\cdot n\,\dd A .
\]
The current basis is normalized so that a port injection column $B_p$ produces
unit terminal current:
\[
  \mathcal I B_p=e_p,
\]
where $\mathcal I$ is the matrix of terminal current functionals.

Port voltage is defined by the work-conjugate functional. The canonical
definition and the exact current-constrained KKT/Schur formulation are given in
Section~\ref{sec:port-schur}. Impressed-current right-hand sides used by an
implementation are admissible only if regression tests show equivalence to that
canonical port definition. The discrete port pair must satisfy the power identity
\[
  \frac12\Rea\{i\adj v\}
  =
  \frac12\Rea\{x_{\rm EM}\adj b(i)\}
\]
up to the independently represented radiation/background channels.

\subsection{Discrete current--charge integral system}

Testing Ohm's law
\[
 \rho_e\mathbf J
 +j\omega\mathbf A
 +\nabla\phi
 =\mathbf E_{\rm impressed}
\]
with the current basis and coupling it to
Equation~\eqref{eq:discrete-continuity} defines the internal mixed operator
\[
 A_{cq}
 =
 \begin{bmatrix}
   R_c+j\omega L_A & -D^\ast \Phi_\rho\\
   D                & j\omega M_\rho
 \end{bmatrix}.
 \label{eq:conductor-block}
\]
For a port-driven problem with no independent impressed field,
\[
  A_{cq}
  \begin{bmatrix}
    c\\q
  \end{bmatrix}
  -
  \mathcal I_{cq}\adj v
  =0,
\]
together with
\[
  \mathcal I_{cq}
  \begin{bmatrix}
    c\\q
  \end{bmatrix}
  =i,
\]
which is the canonical constrained system of
Section~\ref{sec:port-schur}. A lifting implementation may transform this into
a nonzero compatible right-hand side, but must not add that source on top of
the total-current constraint.
With the definition of $D$ above, integration by parts gives the minus
sign in the current--potential coupling. A different signed-divergence
convention may move that sign into $D$, but continuity and the work identity
must be changed consistently.

Here
\[
 (R_c)_{\alpha\beta}
 =
 \int_{\Omega_c}
 \mathbf w_\alpha\adj\rho_e\mathbf w_\beta\,\dd V,
\]
\[
 (L_A)_{\alpha\beta}
 =
 \mu\iint_{\Omega_c\times\Omega_c}
 \mathbf w_\alpha(\mathbf r)\adj
 G_k(\mathbf r,\mathbf r')
 \mathbf w_\beta(\mathbf r')
 \dd V'\dd V,
\]
and $\Phi_\rho$ is the scalar-potential interaction operator.

Equation~\eqref{eq:conductor-block} is the internal current--charge operator;
Section~\ref{sec:port-schur} supplies its canonical work-conjugate port
constraint. Stable implementations may
eliminate charge, use loop/tree or divergence-conforming bases, or use a
regularized mixed formulation. What is invariant is the exact discrete
continuity relation, finite-cross-section current space, and work-conjugate port
pair.

\subsection{Gauge and null-space handling}

The scalar potential is defined up to a constant on electrically isolated
components. One null mode per disconnected charge domain is removed by an
explicit gauge condition, for example
\[
  \mathbf 1^\T M_\rho q=0,
\]
or by a projected solve. Gauge removal must not modify terminal current
constraints.

\subsection{Passivity of the conductor block}

For a passive conductor,
\[
  R_c=R_c\adj\PSD.
\]
The real input power generated by the block system is
\[
  P_c
  =
  \frac12 c\adj R_c c\ge0 .
\]
Thus the conductor contribution to the port impedance is passive before any
neural approximation. Neural corrections must preserve this property by the
structured decoder in Section~\ref{sec:port-macro}.

\subsection{Static and low-frequency limit}

As $\omega\to0$, the formulation must approach the DC conduction/current
constraint without catastrophic cancellation between charge and vector
potential terms. A practical implementation therefore uses frequency-aware
scaling or a loop/tree-like decomposition. The release suite must contain a
low-frequency sequence specifically checking this asymptotic behavior.

\subsection{Relation to PEEC}

After basis testing, the resistance, vector-potential, and scalar-potential
blocks have the same physical interpretation as partial elements in PEEC
\cite{ruehli1974peec}. vNext uses this circuit interpretation while retaining
continuous superelliptic geometry and higher-order finite-cross-section modes.
